% page 342 du fac-simile (image p-341 du scan)
% bloc 157.5 x 246.3 mm ; marge gauche 8.0 mm
\pgc{-12.700mm}{0.000mm}{
\l{\cel{3}334}
\l{}
\l{}
\l{\sou{lektoro}. I. Sorto di profesoro adyunta en ula universitati.}
\l{\cel{3}- II. (\sou{liturgio katolika}). Kleriko kun la duesma ek la}
\l{\cel{3}qua ordeni minora, qua havis olim kom tasko lektar, kan-}
\l{\cel{3}tar la lecioni, epistoli, e c., en la ceremonii kultala.}
\l{\cel{3}- DEFIS.}
\l{}
\l{\sou{lemo.}\cel{2}(\sou{matem.)}\cel{2}Propoziciono preliminara quan onu establisas}
\l{\cel{3}por preparar la demonstro di altra propoziciono. - DEFIRS.}
\l{}
\l{\sou{lemniskato.}\cel{1}(\sou{geom.}) Kurvo ek la grado quaresma, qua havas}
\l{\cel{3}la formo di la cifro \textquotedbl{}8\textquotedbl{}, e qua esas la ligilo di omna}
\l{\cel{3}punti tale ke la produto di lia disti a du punti fixa}
\l{\cel{3}esas konstanta. - DEFI.}
\l{}
\l{\sou{lemuriano}. (\sou{zool.}) Familio de simii, di qui la tipo esas}
\l{\cel{3}la makio. - DEFIS.}
\l{}
\l{\sou{lengo}. (\sou{zool.}) Speco di gado, di qua la korpo esas oblonga.}
\l{\cel{3}- L. gadus molva.}
\l{}
\l{\sou{lenio}. Peco di kalorizo-ligno, segita o tranchita tale ke}
\l{\cel{3}onu povas pozar lu aden kameno, aden furnelo. - IS.}
\l{}
\l{\sou{lenso}.- I. (\sou{bot.}) Grano di leguminoso quan onu manjas kom}
\l{\cel{3}legumo pos lua sikigeso. - II. (\sou{fiziko}).Vitro-bloko,}
\l{\cel{3}taliita lenso-forme, finanta sive per du surfaci sferala,}
\l{\cel{3}sive per surfaco plana, qua produktas deviaco regulala}
\l{\cel{3}di la lumo-radii. - EF.}
\l{}
\l{\sou{lenta}. I. Qua ne iras rapide. - II. Qua trodurigas to quon}
\l{\cel{3}lu esas aganta o facanta. - FIS.}
\l{}
\l{\sou{lentigino.}\cel{2}(\sou{che ula personi}). Altereso naturala, grano-}
\l{\cel{3}forma, di la koloro di lia pelo. - EFIS.}
\l{}
\l{\sou{lentisko}. (\sou{bot.)}\cel{2}Arboreto qua kreskas en la Franca provinco}
\l{\cel{3}Provence ed en Italia, e di qua un speco donas la suko}
\l{\cel{3}rezinoza de qua onu facas liquoro e mastico. L.\cel{1}\sou{Pistacia}}
\l{\cel{3}\sou{lentiscus.}\cel{1}- EFIS.}
\l{}
\l{\sou{leono.}\cel{1}I(zool.) Mamifero karnivora, ek la familio \textquotedbl{}felini\textquotedbl{},}
\l{\cel{3}kun pilkoloro flava, kun la shultri kovrata (che la masku-}
\l{\cel{3}lo) da dika krinaro. - II. (\sou{astrol}.) Signo di zodiako,}
\l{\cel{3}reprezentata da leono.- III. (\sou{metaf.}) (l) viro brava,}
\l{\cel{3}audacoza quale leon(in)o. (2) Personego tre segun-moda, e}
\l{\cel{3}tre admirata, che la mondumani. - EFIS.}
\l{}
\l{\sou{leopardo.}\cel{2}(\sou{zool.)}\cel{2}Animalo karnavida, ek la genero \textquotedbl{}kato\textquotedbl{},}
\l{\cel{3}kun pelo makuletoza. - L.\cel{1}\sou{felis leopardus.}\cel{1}- DEFIRS.}
\l{}
\l{\sou{lepidoptero}. (\sou{zool.}) Familio ek insekti kun ali squamoza, no-}
\l{\cel{3}mata vulgare \textquotedbl{}papilioni\textquotedbl{}, qui subisas metamorfosi kompleta,}
\l{\cel{3}e di qui onu nomizas la larvo \textquotedbl{}raupo\textquotedbl{}, e la nimfo \textquotedbl{}kriza-}
\l{\cel{3}lido\textquotedbl{}. - DEFIS.}
}
